\chapter{Thất bại của mô hình, khủng hoảng và nhu cầu về dữ liệu}
\label{ch:M6L4LessonNotes}

\section*{Lesson Notes}

MODULE 6 \textbar{} BÀI 4

{\small
\begin{longtable}[]{@{}ll@{}}
\toprule\noalign{}
\endhead
\bottomrule\noalign{}
\endlastfoot
\textbf{Thời gian đọc} & 30 phút \\
\textbf{Kiến thức nền} & Khóa Thị trường tài chính \\
\end{longtable}
}

\textbf{Từ khóa:} Mô hình mẫu (model paradigm), phản tác dụng (backfire)

\emph{Trong bài học cuối của học phần này, ta tóm lược những bài học chính về dữ liệu và vai trò của dữ liệu trong việc giải quyết các thách thức tài chính đã bàn ở khóa Thị trường tài chính. Cụ thể, bài học đề cập đến thất bại của mô hình và khủng hoảng, qua đó cho thấy vì sao cần có dữ liệu.}

Đôi khi mô hình thất bại. Vấn đề không nằm ở thị trường mà ở lăng kính biểu diễn đơn giản hóa, qua đó ta tóm lược, giản lược hoặc hệ thống hóa thế giới. Những mô hình đã hiệu quả hàng thập kỷ, được kiểm định và tái lập vẫn có thể thất bại. Những mô hình được chứng minh bằng suy diễn toán học và đăng trên tạp chí bình duyệt vẫn có thể thất bại. Những mô hình nổi tiếng, được trích dẫn nhiều và mang về giải Nobel cho tác giả vẫn có thể thất bại. Những mô hình có mặt trong phần mềm dùng khắp thế giới vẫn có thể thất bại. Vì sao lại như vậy?

Có đủ loại mô hình: mô hình dự báo biến động và quyền chọn, mô hình xác định phần bù rủi ro; mô hình giao dịch theo giá trị, ít giao dịch; mô hình giao dịch tần suất cao hoạt động trong tích tắc; mô hình cấu trúc vi mô thị trường mô tả chi tiết các lệnh trên sàn; mô hình kinh tế vĩ mô toàn cầu về GDP, việc làm, lạm phát và chính sách tiền tệ. Mô hình có đủ hình dạng và quy mô.

Điểm chung của mọi mô hình là thường giúp thấu hiểu các thách thức đã bàn: rủi ro tín dụng và tài trợ, biến động, tương quan, đòn bẩy, phi tuyến, thanh khoản và quy định. Tuy nhiên, khi mô hình thất bại, kết quả và cách diễn giải của chúng có thể gây hiểu lầm.

Hãy xem xét mười bài học rút ra từ thất bại của mô hình.

\section{Tài chính không phải là vật lý}

Các định luật vật lý không phải là định luật của tài chính. Cần phân biệt rõ điều đúng trong các mô hình lý thuyết, duy lý với điều đúng trong thực tế. Tài chính không phải là vật lý: mô hình tài chính phản ánh hành vi con người, vốn là sự pha trộn giữa duy lý và phi lý. Vì vậy, mô hình tài chính có thể thất bại nếu giả định mọi người đều hành xử duy lý. Đôi khi thị trường dao động giữa trạng thái hiệu quả cao và trạng thái bị chi phối bởi các thiên lệch nhận thức (bias) cực đoan. (Xem Giả thuyết Thị trường Thích ứng (Adaptive Market Hypothesis) của Andrew Lo \url{https://web.mit.edu/Alo/www/Papers/JPM2004_Pub.pdf} để có thảo luận thú vị về sự dao động này).

Yếu tố phi lý mới xuất hiện trong kinh tế học gần đây hơn nhiều so với các giả định duy lý. Chỉ trong 20 năm qua, giải Nobel mới được trao cho các nhà nghiên cứu kinh tế học hành vi (ví dụ Kahneman, Thaler). Tài chính hành vi nhắc rằng nhiều giả định nền tảng của tài chính thực tế là sai sót. Do đó, tài chính không có những định luật bất biến như các ngành khoa học tự nhiên. Một số ý tưởng này được nêu trong một tuyên ngôn viết sau Đại khủng hoảng tài chính (Wilmott), nhắc rằng thiên lệch của con người trong ra quyết định khiến mọi "định luật tài chính" không thể tái lập.

Hãy xét đợt tăng giá gần đây của cổ phiếu Game Stop. Trong nhiều bài học rút ra, có việc con người có thể thao túng thị trường chứng khoán vì những lý do vượt ngoài đặc điểm rủi ro và lợi suất. Xem và đọc bài của Zhao (\url{https://www.atlantis-press.com/proceedings/icemci-21/125966071}).

\section{Các vị thế dùng đòn bẩy có thể khuếch đại lợi nhuận hoặc làm trầm trọng thêm khoản lỗ.}

Giả sử bạn có một mô hình đã được kiểm định tốt và khai thác được alpha. Bạn có thể bị thúc đẩy dùng đòn bẩy để thu về lượng alpha lớn hơn, chẳng hạn đầu tư gấp 8 lần vào chiến lược với kỳ vọng doanh thu ròng cũng gấp 8 lần. Đòn bẩy cho phép bỏ ra ít vốn hơn, vay thêm và tạo lợi suất lớn hơn. Quả thực chiến lược này có thể mang lại lợi suất cao hơn, nhưng cũng có thể gây lỗ lớn hơn. Điều chắc chắn là nó làm rủi ro tăng lên: đòn bẩy càng nhiều thì độ bất định quanh giá trị kỳ vọng càng lớn. Vị thế dùng đòn bẩy không phân biệt hướng đi: nó có thể tăng lãi hoặc làm khoản lỗ sâu hơn.

Để xem các ví dụ chi tiết hơn, hãy đọc bài "Leverage and Risk Taking under Moral Hazard" của Hott (\url{https://www.google.com/url?q=https\%3A\%2F\%2Flink.springer.com\%2Farticle\%2F10.1007\%2Fs10693-021-00359-8}).

\section{Dễ nhầm biến với tham số: biến thay đổi, còn tham số là hằng số.}

Làm sao biết đại lượng nào giữ nguyên và đại lượng nào thay đổi trong một mô hình? Khi mô hình hóa một hệ thống, ta chọn một số đại lượng làm tham số và số khác làm biến. Tham số là các hằng số được ước lượng từ dữ liệu lịch sử hoặc suy ra từ giá thị trường. Biến là các đại lượng động, thay đổi theo thời gian. Rất dễ coi biến là tham số, và trong bối cảnh tài chính, người ta thường gọi nhầm một biến là tham số.

Thiếu nền tảng toán và thống kê vững, ta dễ giả định nhiều thứ là hằng số: biến động, lãi suất, tương quan, thậm chí cả phân phối. Bộ công cụ phân tích càng mạnh thì ta càng có thể nới lỏng giả định và nâng cấp mô hình, chẳng hạn cho biến động hay lãi suất là biến thay vì tham số. Trớ trêu là đôi khi ta lại muốn một số đại lượng là tham số thay vì biến. Người không chuyên chỉ có thể dựa vào các mô hình do người khác xây dựng. Người có chuyên môn kỹ thuật không chỉ dùng mà còn điều chỉnh được mô hình cho mục đích riêng bằng cách nới lỏng giả định, dùng thuật toán mạnh hơn và các phương pháp hiệu chỉnh, tối ưu hóa nâng cao. Hơn nữa, nếu bổ sung sự sáng tạo, tư duy phản biện và chuyên môn lĩnh vực, họ có thể tạo ra các mô hình tùy biến.

Điểm cuối này cho thấy vì sao cần một chương trình thạc sĩ, chứ không chỉ một khóa học hay một video về "cách trở thành kỹ sư tài chính". Tài liệu cung cấp lý thuyết và cách triển khai các công cụ, sau đó dùng chúng để minh họa việc giải quyết vấn đề trong kỹ thuật tài chính. Nếu mô hình biến động hằng số tỏ ra không đủ (như khi cố tái tạo nụ cười biến động), ta có thể nâng cấp lên mô hình biến động ngẫu nhiên hai nhân tố, đòi hỏi hiểu biết toán, hiệu chỉnh thống kê và triển khai tính toán. Điều tương tự áp dụng cho lãi suất động, tương quan động hay bất kỳ đại lượng nào ta muốn biểu diễn là biến thiên theo thời gian. Cái giá của mô hình tinh vi hơn là nền tảng để hiểu chúng, độ phức tạp tính toán khi triển khai, và sự thấu hiểu để biết khi nào chúng hiệu quả và khi nào không.

\section{Ta có thể dùng mô hình sai cách, do không hiểu hướng dẫn hoặc giới hạn của chúng.}

Mô hình có thể vận hành đúng, nhưng ta lại dùng chưa đúng: đánh giá thấp hoặc cao một tham số, làm sai hướng dẫn, không hiểu các giả định mô hình áp đặt, hoặc dùng mô hình trong điều kiện nó không lường trước. Lỗi không nằm ở mô hình mà ở cách ta sử dụng.

Ví dụ đơn giản là sự cẩu thả. Khi gửi lệnh lên sàn, lệnh là một "mô hình": ta chỉ định giá và khối lượng. Giả sử muốn bán 1 đơn vị với giá 610.000 yên, nhưng lại nhầm giá với khối lượng và gửi lệnh bán 610.000 đơn vị với giá 1 yên. Nghe khó tin? Việc này đã xảy ra vào tháng 12 năm 2005. Lệnh rất bất thường: giá không phải 1 yên mà là 610.000 yên (khoảng \$5.041), và số cổ phiếu trong lệnh gấp hơn 40 lần số công ty đang phát hành. Dù vậy, sàn vẫn không từ chối lệnh. Thiệt hại gần một phần tư tỷ đô la. Lỗi này thuộc về rủi ro vận hành. Phần mềm được đưa vào sản xuất mà thiếu kiểm thử và thẩm định như thế nào? Hãy hình dung một mô hình có hàng nghìn bộ phận và việc đánh giá xem nó có sai sót cơ bản hay không.

Khó hơn là khi các giả định hoặc điều kiện xây dựng mô hình không được tuân thủ. Có những ví dụ bi thảm: năm 1986, tàu con thoi Challenger nổ tung, giết chết cả bảy phi hành đoàn. Thủ phạm là một bộ phận, vòng đệm O-ring. Richard Feynman, một trong những nhà vật lý xuất sắc nhất thế giới, đã giải thích nguyên nhân. Khi chờ phát biểu trước hội đồng điều tra, ông bỏ các vòng O-ring vào nước đá; đến lượt mình, ông lấy ra và cho thấy chúng mất độ đàn hồi như giả định. Các vòng này chỉ được thử nghiệm ở Florida trong thời tiết ấm, cao hơn nhiều so với điểm đóng băng, nhưng vào ngày phóng tháng 1, Florida lạnh bất thường, ở mức nhiệt chưa từng được thử nghiệm.

Đây là ví dụ áp dụng kết quả mô hình ở mức nhiệt mà nó hầu như chưa được kiểm tra; theo thuật ngữ thống kê, đó là ngoại suy (extrapolation). Việc thử nghiệm đã không xét đến hiệu quả của O-ring trong môi trường lạnh hơn bình thường 20 độ F (11 độ C). Dữ liệu và câu chuyện được trình bày trong bài của Wujek (\url{https://www.google.com/url?q=https\%3A\%2F\%2Fonlineethics.org\%2Fcases\%2Fnumerical-design-problems-ethical-content\%2Fchallenger-o-ring-data-analysis}). Hãy đọc bài này. Mô hình thất bại có thể rất tốn kém, đánh đổi bằng tính mạng và tài sản.

\section{Tính thanh khoản có thể khiến mô hình cho kết quả ngoài dự kiến.}

Giả sử bạn có một giao dịch rất tốt và muốn mở rộng quy mô bằng cách mua thêm. Nếu vị thế quá lớn, thanh khoản của chứng khoán bắt đầu bị đặt dấu hỏi: liệu có tác động lên thị trường (market impact) hay không? Một giao dịch có lãi vẫn có thể trở thành thua lỗ nếu mở rộng quá mức, vì tác động thị trường triệt tiêu alpha của chiến lược. Mỗi nhà giao dịch đều có một mô hình khai thác alpha từ thị trường. Khi giao dịch ngày càng nhiều, họ gặp trượt giá (slippage): tác động thị trường làm tăng giá khi vào lệnh và hạ giá khi thoát lệnh. Liệu có một mô hình tương ứng cho tác động thị trường của chính mình? Dù mua hay bán, tác động thị trường làm giá dịch chuyển theo hướng ăn mòn alpha của chiến lược. Nếu không kiểm soát, trượt giá có thể lấn át mọi lợi thế. Chúng ta sẽ bàn về trường hợp Zillow ngay sau đây.

\section{Sự tồn tại của một mô hình, hay niềm tin của tác giả vào nó, không có nghĩa là mô hình hiệu quả trong thực tế.}

"Chúng tôi dùng mô hình này từ xưa đến nay." "Đây là mô hình ai cũng dùng." "Mô hình này có tuổi đời lâu hơn cả tuổi của bạn!" "Tại sao tôi phải tin mô hình của bạn khi những người nổi tiếng và hàng tỷ đô la đã dùng mô hình này?"

Bạn được giao một mô hình và thường được bảo rằng nó hiệu quả, điều này tạo cảm giác an toàn trọn vẹn. Bạn sẽ gặp tình huống này nhiều lần trong chương trình. Chẳng hạn, trong Định giá phái sinh (Derivative Pricing), bạn định giá một quyền chọn bằng mô hình đã có 50 năm tuổi; bạn tính giá và các độ nhạy nhưng ban đầu không đặt câu hỏi nhiều. Dần dần, bạn nhận ra một số giả định của mô hình có sai sót, ví dụ được bảo biến động là hằng số, nhưng sau đó thấy thị trường quyền chọn không quan sát được điều đó. Vì sao người ta bám lấy mô hình dù thấy bằng chứng ủng hộ rất ít? Chúng ta sẽ bàn về các thiên lệch (bias) nhận thức và hành vi trong Quản lý danh mục đầu tư (Portfolio Management); chẳng hạn, thiên lệch quá tự tin có thể che mắt một người trước sự thật vì học vấn, kinh nghiệm và vị trí của họ. Thị trường có thể khiến ta khiêm nhường.

Mô hình lâu đời, từng đoạt giải Nobel, phổ quát, dễ triển khai, được trích dẫn thường xuyên, cực kỳ phổ biến hay trực quan không nhất thiết là mô hình tốt nhất để dùng. Động lực thị trường thay đổi. Hãy xét các mô hình cho cổ phiếu: cổ phiếu chịu tác động từ quyết định của ngân hàng trung ương; việc thêm hoặc loại bỏ cổ phiếu khỏi các chỉ số lớn có thể ảnh hưởng đến chúng; ngày đáo hạn quyền chọn tác động lên thị trường cổ phiếu cơ sở. Ngay cả sự tồn tại và mở rộng của tiền mã hóa cũng ảnh hưởng đến cách thị trường cổ phiếu giao dịch! Các mô hình định giá vốn hiệu quả trong giao dịch cổ phiếu có thể bị gián đoạn tạm thời (nếu chưa phải đổ vỡ) bởi sự mở rộng và sụp đổ của các sàn giao dịch tiền mã hóa. Với vai trò kỹ sư, bạn dùng dữ liệu và năng lực mô hình hóa để xác định độ nhạy, các điểm gãy cấu trúc hay sự chuyển đổi chế độ. Trớ trêu thay, đôi khi bạn lại dùng thêm nhiều mô hình để đánh giá hiệu quả xem các mô hình đang có có phù hợp hay không. Trong một số tình huống khắc nghiệt, bạn có thể cần trao đổi với nhà nghiên cứu, nhà quản lý danh mục đầu tư và nhà quản lý rủi ro để xác định mô hình có thiếu độ tin cậy hay thiếu kiểm định hay không. Phụ thuộc mù quáng vào các mô hình chưa được kiểm chứng có thể tạo cảm giác an toàn giả tạo.

Hầu như không ai nhận được thông báo: "Mô hình này không còn hiệu quả. Hãy dùng mô hình khác." Mô hình có thể được dùng rất lâu sau "ngày hết hạn". Hãy học cách tư duy phản biện về mô hình, và đừng biện minh cho việc sử dụng bằng ngụy biện viện dẫn quyền uy.

\section{Các mô hình giả định tuyến tính có thể thất bại do phi tuyến tính, đặc biệt khi có đòn bẩy.}

Ta thường giả định mô hình có tính tuyến tính: tăng đầu vào thì đầu ra tăng theo một mức đã biết. Giả sử tỷ lệ là 1,5. Khi đó, nếu đầu vào tăng gấp đôi thì đầu ra tăng gấp ba. Trong nhiều trường hợp khác, ta lại thấy có nhiều phi tuyến tính. Như đã thấy, tính lồi (convexity) khiến thay đổi nhỏ chỉ tạo kết quả nhỏ, nhưng thay đổi lớn tạo kết quả lớn hơn một cách không tương xứng.

Tương tự, tính phi tuyến từng khúc trong payoff của quyền chọn có thể đưa quyền chọn vào trạng thái có lời (in the money) hoặc không có lời (out of the money). Việc mức tăng đưa bạn sang phía nào của giá thực hiện (strike) là rất quan trọng! Tác động thị trường (market impact) là ví dụ về độ nhạy có thể ban đầu tuyến tính (với lượng giao dịch tương đối nhỏ) rồi trở nên rất phi tuyến với lượng giao dịch lớn hơn. Sai lầm thường là đánh giá thấp nghiêm trọng một chi phí, tổn thất hoặc rủi ro, có thể gây ra khủng hoảng.

\section{Mô hình được xây dựng trên các giả định, vì vậy hãy nắm rõ các điều kiện.}

Mọi mô hình đều có giả định. Nhiều người không nhận ra các giả định này vì cần hiểu sâu hơn việc chỉ chạy code để dùng mô hình. Có những điều kiện quyết định có nên dùng mô hình hay không, chẳng hạn: Gaussian so với phi Gaussian, dừng so với không dừng, động lượng (momentum) so với hồi quy về trung bình, không âm so với có thể âm. Giả định thường quyết định mô hình thành hay bại, và đôi khi điều kiện không thỏa mãn khiến mô hình mất hiệu lực. Lỗi phổ biến là chạy hồi quy trên các quá trình không dừng rồi thấy một thống kê kiểm định có vẻ có ý nghĩa, trong khi đó có thể chỉ là hồi quy giả (spurious regression). Đây là lỗi của người vận hành hơn là lỗi của mô hình, vì ngay từ đầu mô hình đã không phù hợp để sử dụng. Do đó, việc xây dựng nền tảng vững chắc là thiết yếu. Nhiều sinh viên và người đi làm chỉ chạy mô hình theo thủ tục mà không đánh giá xem mô hình có phù hợp hay không. Có thể khắc phục dễ dàng bằng cách nghiên cứu các tính chất toán học nền tảng của mô hình.

Các điều kiện khác đơn giản liên quan đến giới hạn tự nhiên của biến mục tiêu. Mô hình dự báo giá cổ phiếu âm là đáng ngờ, vì giá cổ phiếu không thể âm. Tuy vậy, mô hình cấm lãi suất âm cũng đáng ngờ không kém, vì lãi suất âm có tồn tại. Tương tự, giá hợp đồng tương lai dầu có thể âm. Không phải lúc nào cũng dễ xác định đại lượng nào buộc phải không âm.

\section{Mọi mô hình đều là một mô thức (paradigm) cần được kiểm định để đánh giá đúng.}

Mô hình phản ánh các mô thức. Rất thường xuyên, bạn nhìn thế giới thông qua chính mô hình mình đang dùng. Chẳng hạn, nếu tin rằng lợi suất của một cổ phiếu tuân theo phân phối Gauss, bạn sẽ gọi một khoản lỗ rất lớn là biến động 5 sigma, một sự kiện thực sự rất hiếm. Nhưng nếu phân phối không thực sự là Gauss thì sao? Các sự kiện 5 sigma là cực kỳ khó xảy ra. Nhiều phân phối lợi suất có độ nhọn dư (excess kurtosis), khiến chúng không phải Gauss. Vì vậy, ta nên bỏ "cặp kính Gauss" và đeo cặp kính nhìn thế giới qua một phân phối phù hợp hơn.

Mô hình cần được kiểm định. Rất khó biết mô hình của bạn sai nếu bạn nhìn thế giới như thể các giả định mặc nhiên đúng. Các ngân hàng được phép dùng mô hình nội bộ để chạy kiểm tra sức chịu đựng (stress test) đối với rủi ro tín dụng, nhưng các mô hình này cần được kiểm định kỹ lưỡng. Với các ngân hàng trung ương hay quy định Basel, họ yêu cầu những tiêu chí nhất định để kiểm định mô hình.

Là học viên của chương trình, bạn sẽ học các mô hình nền tảng rồi bổ sung độ phức tạp cho chúng. Chẳng hạn, trong lộ trình Học máy, bạn ôn lại bình phương tối thiểu rồi học nhiều mở rộng. Có hai nhận định dẫn dắt quá trình này. Thứ nhất, học mô hình đơn giản giúp xây nền tảng để học mô hình nâng cao. Thứ hai, nếu mô hình đơn giản hoạt động tốt thì ta có thể dùng chúng; nhiều khi mô hình đơn giản lại tốt nhất! Tuy nhiên, nếu các mô hình này khớp thiếu (underfit), ta còn các phương pháp khác để dùng. Hiểu "tại sao" và "như thế nào" đằng sau các mô hình sẽ mài giũa bộ kỹ năng của bạn.

\section{Mô hình có thể phản tác dụng.}

Mô hình có thể phản tác dụng: mô hình nhằm giảm rủi ro có thể làm rủi ro tăng, mô hình nhằm tăng lợi suất thực tế lại làm giảm lợi suất. Ngoài rủi ro, mô hình còn chứa bất định: ta có thể không biết phân phối, thậm chí không biết các kết cục. Hãy xét một ví dụ thực tế.

Gần đây, Zillow, nhà cung cấp phân tích thông tin bất động sản, triển khai một chiến lược tự thân và bị phản tác dụng. Họ quyết định mua đi bán lại (flip) nhà ở, vì sở hữu lượng dữ liệu khổng lồ và tin rằng dữ liệu này tạo lợi thế trên thị trường do ít công ty khác có được. Dữ liệu không chỉ gồm giá niêm yết và giá chốt, mà còn lượt truy cập, tìm kiếm và dữ liệu khác từ cổng tìm kiếm rất phổ biến của họ. Chẳng hạn, khu vực hay con phố có nhiều lượt tìm kiếm trở thành bất động sản được ưa chuộng. Zillow tin mình nắm rõ bất động sản nào được ưa chuộng và sẽ bán bằng hoặc cao hơn giá niêm yết. Ngoài ra, họ còn có thuật toán độc quyền định giá nhà, gọi là Zestimate.

Zillow mua gom quyết liệt để bán lại, đồng thời thao túng thị trường nhẹ bằng cách nâng giá dần khi mua, hy vọng tạo xu hướng khiến toàn bộ số nhà đó được định giá cao hơn. Tiếc là thị trường nhận ra chiến lược này, tạo cảm xúc rất xấu rằng họ đang thao túng giá nhà. Họ cũng đánh giá thấp việc thiếu đa dạng hóa: các vụ mua tập trung quanh những điểm nóng. Họ cũng đánh giá thấp rủi ro thanh khoản: việc cải tạo và bán nhà mất nhiều thời gian hơn dự kiến. Họ lại mua khá muộn trong giai đoạn bong bóng, trả giá gần đỉnh thị trường. Khi bán hàng tồn kho, giá bán trung bình thấp hơn khoảng 80.000 đô la mỗi căn. Tổng cộng họ lỗ hàng trăm triệu đô la, phải sa thải nhân viên và cuối cùng rút khỏi mảng kinh doanh đầu tư.

Dù có dữ liệu, khả năng tiếp cận, kinh nghiệm và các mô hình định giá, họ vẫn không tính đến những vấn đề then chốt như tương quan, thanh khoản và danh tiếng. Lợi thế họ tin mình có rốt cuộc phản tác dụng thành một khoản đầu tư thảm họa.

Mô hình có thể phản tác dụng vì không dễ biết được ảnh hưởng của chính ta lên thị trường.

\section{Kết luận}

Trong các dự án capstone, sinh viên thường chạy kiểm định lùi (backtesting) với giả định có thể mua bán tại giá đóng cửa. Họ dùng lượng vốn rất lớn và báo cáo kết quả phi thường, không có tác động thị trường hay chênh lệch giá mua-bán (bid-ask spread). Họ cho rằng kết quả backtest trên giấy phản ánh chính xác điều sẽ xảy ra. Nhưng khi tham gia thị trường, họ chắc chắn sẽ gặp chênh lệch giá mua-bán và có thể chịu tác động thị trường, tùy quy mô giao dịch.

Như trước đây, điểm chung của mọi mô hình là cung cấp cái nhìn sâu sắc về các thách thức đã bàn: rủi ro tín dụng và tài trợ, biến động, tương quan, đòn bẩy, phi tuyến, thanh khoản và quy định. Tuy nhiên, khi bị dùng sai, kết quả và diễn giải của chúng tốt nhất là vô nghĩa (tức không dùng được), tệ nhất là gây hiểu lầm nghiêm trọng. Cảm giác an toàn giả tạo đó có thể dẫn đến rủi ro hệ thống lan rộng.

Trong bài học cuối này, ta thấy mô hình thất bại vì ba nguyên nhân bao trùm:

\begin{enumerate}
\def\labelenumi{\arabic{enumi}.}
\item
  chúng có thể được xác định sai,
\item
  chúng có thể bị sử dụng sai, và
\item
  chúng có thể bị sử dụng không phù hợp.
\end{enumerate}

Với những hiểu biết mới này, ta tiếp tục. Ta sẽ đón nhận các thách thức tài chính này và tận dụng kinh nghiệm phân tích thực nghiệm bằng cách dùng dữ liệu trong nhiều mô hình, dẫn ta qua hành trình nhiều học kỳ gồm kinh tế lượng, phân tích chuỗi thời gian, quá trình ngẫu nhiên, học máy, học sâu và hơn thế nữa. Dù dùng mô hình nào, hãy luôn thận trọng về cách sử dụng (và cố gắng không lạm dụng) chúng.

\begin{center}\rule{0.5\linewidth}{0.5pt}\end{center}

Copyright 2024 WorldQuant University. This content is licensed solely for personal use. Redistribution or publication of this material is strictly prohibited.
